%% ma: L12-B08-0147
%% lop: 12
%% chuong: 2
%% bai: 8
%% dang: 12.8.49
%% loai: TN4
%% muc: VD
%% dap_an: "A"
%% nguon: "CĐ2.3.pdf giáo viên gửi 27/10/2026 (M07, Lê Minh Khôi – Nguyễn Anh Khôi, Chương II Chủ đề 3 Biểu thức tọa độ của các phép toán vectơ), Câu 63"
%% kien_thuc: "Bài 8 (điều kiện vuông góc, biểu thức tọa độ)"
%% trang_thai: chinh-thuc
%% ghi_chu: "Giáo viên đã duyệt 27/10/2026"
\begin{ex}
Trong không gian $Oxyz$, cho tam giác $JKQ$ có tọa độ đỉnh $J(1;1;0)$, $K(0;2;1)$. Biết rằng $H(-1;2;1)$ là trực tâm của tam giác $JKQ$. Tọa độ của điểm $Q$ là
\choice{\True $Q(1;3;2)$}{$Q(-1;3;2)$}{$Q(3;1;2)$}{$Q(1;-3;-2)$}
\loigiai{Thử các phương án với ba điều kiện $\overrightarrow{JH}\cdot\overrightarrow{KQ}=0$, $\overrightarrow{KH}\cdot\overrightarrow{JQ}=0$ và $H$ thuộc mặt phẳng $(JKQ)$ (ba vectơ $\overrightarrow{JK}$, $\overrightarrow{JQ}$, $\overrightarrow{JH}$ đồng phẳng).

Với $Q(1;3;2)$: $\overrightarrow{JH}\cdot\overrightarrow{KQ}=(-2;1;1)\cdot(1;1;1)=0$, $\overrightarrow{KH}\cdot\overrightarrow{JQ}=(-1;0;0)\cdot(0;2;2)=0$ và $\overrightarrow{JH}=2\overrightarrow{JK}-\dfrac12\overrightarrow{JQ}$. Đúng; các phương án khác không thỏa.}
\end{ex}
